% GENERATED FILE. Edit canonical lessons, then run npm run generate:books.
% canonical-source-hash: fnv1a64:e4391ffc77b7aa4c
% canonical-lessons: FR-C472-nostalgie, FR-C472-regretter, FR-C472-melancolie, FR-C472-emouvoir, FR-C472-bouleverser

\chapter{Souvenirs et regrets}
\label{ch:fr-souvenirs-et-regrets}
\hlchaptermodality{472}

\begin{chapteropening}
\textbf{By the end of this chapter:} \emph{I can talk about memories, regrets and being moved.}

Complete the last lesson of chapter 472: I can talk about memories, regrets and being moved.
\end{chapteropening}

\section[la nostalgie]{la nostalgie — nostalgia, longing for the past}

\label{lesson:FR-C472-nostalgie}



\begin{quote}
Before the new one: say the French for to sob, then the French for to comfort.
\end{quote}

\subsection*{You'll want to know: la nostalgie}
\textbf{la nostalgie} — \textquotedblleft{}nostalgia, longing for the past\textquotedblright{}.

\begin{quote}
Je pense à mon enfance avec nostalgie.
\end{quote}

\begin{quote}
\emph{I think of my childhood with nostalgia.}
\end{quote}

\emph{Avoir la nostalgie de} — to miss, to long for.

\begin{cousinweb}
\emph{la nostalgie} comes from New Latin \emph{nostalgia} (EtymDB), coined from Greek \emph{nóstos} \textquotedblleft{}return home\textquotedblright{} and \emph{álgos} \textquotedblleft{}pain\textquotedblright{}: it first named homesickness.
\end{cousinweb}

\begin{grammarlens}[title={What you've built}]
One more word for telling, explaining and describing at B1.
\end{grammarlens}

\subsection*{Your turn}
\begin{itemize}
\raggedright
  \item \emph{Say it:} \emph{la nostalgie}
  \item \emph{Say it:} \emph{la nostalgie}, once more
  \item \emph{Say it:} a sentence of your own with \emph{la nostalgie}
  \item \emph{Recall:} say the French for to sob, then the French for to comfort, then say \emph{la nostalgie} again
\end{itemize}

\begin{tcolorbox}[breakable,colback=teal!4,colframe=teal!35!black,title={Before you move on}]
What does la nostalgie mean? (\textquotedblleft{}Nostalgia, longing for the past\textquotedblright{}.) Say it once more.
\end{tcolorbox}
\section[regretter]{regretter — to regret; to miss}

\label{lesson:FR-C472-regretter}



\begin{quote}
Before the new one: say the French for to comfort, then the French for nostalgia.
\end{quote}

\subsection*{You'll want to know: regretter}
\textbf{regretter} — \textquotedblleft{}to regret; to miss\textquotedblright{}.

\begin{quote}
Je regrette mon ancien quartier.
\end{quote}

\begin{quote}
\emph{I miss my old neighbourhood.}
\end{quote}

\emph{Je regrette} also means I am sorry: \emph{je regrette, c'est complet}.

\begin{cousinweb}
\emph{regretter} comes from Old French \emph{regreter} \textquotedblleft{}to lament the dead\textquotedblright{} (EtymDB), built from \emph{re-} and an old verb for weeping, probably of Germanic origin.
\end{cousinweb}

\begin{grammarlens}[title={What you've built}]
One more word for telling, explaining and describing at B1.
\end{grammarlens}

\subsection*{Your turn}
\begin{itemize}
\raggedright
  \item \emph{Say it:} \emph{regretter}
  \item \emph{Say it:} \emph{regretter}, once more
  \item \emph{Say it:} a sentence of your own with \emph{regretter}
  \item \emph{Recall:} say the French for to comfort, then the French for nostalgia, then say \emph{regretter} again
\end{itemize}

\begin{tcolorbox}[breakable,colback=teal!4,colframe=teal!35!black,title={Before you move on}]
What does regretter mean? (\textquotedblleft{}To regret; to miss\textquotedblright{}.) Say it once more.
\end{tcolorbox}
\section[la mélancolie]{la mélancolie — melancholy, gloom}

\label{lesson:FR-C472-melancolie}



\begin{quote}
Before the new one: say the French for nostalgia, then the French for to regret.
\end{quote}

\subsection*{You'll want to know: la mélancolie}
\textbf{la mélancolie} — \textquotedblleft{}melancholy, gloom\textquotedblright{}.

\begin{quote}
L'automne me donne toujours un peu de mélancolie.
\end{quote}

\begin{quote}
\emph{Autumn always makes me a little melancholy.}
\end{quote}

The adjective is \emph{mélancolique}.

\begin{cousinweb}
\emph{la mélancolie} comes through Latin \emph{melancholia} from Greek \emph{melankholía} (EtymDB), built on \emph{mélas} \textquotedblleft{}black\textquotedblright{} and \emph{khol\'{\={e}}} \textquotedblleft{}bile\textquotedblright{}: old medicine blamed sadness on black bile.
\end{cousinweb}

\begin{grammarlens}[title={What you've built}]
One more word for telling, explaining and describing at B1.
\end{grammarlens}

\subsection*{Your turn}
\begin{itemize}
\raggedright
  \item \emph{Say it:} \emph{la mélancolie}
  \item \emph{Say it:} \emph{la mélancolie}, once more
  \item \emph{Say it:} a sentence of your own with \emph{la mélancolie}
  \item \emph{Recall:} say the French for nostalgia, then the French for to regret, then say \emph{la mélancolie} again
\end{itemize}

\begin{tcolorbox}[breakable,colback=teal!4,colframe=teal!35!black,title={Before you move on}]
What does la mélancolie mean? (\textquotedblleft{}Melancholy, gloom\textquotedblright{}.) Say it once more.
\end{tcolorbox}
\section[émouvoir]{émouvoir — to move, to touch (emotionally)}

\label{lesson:FR-C472-emouvoir}



\begin{quote}
Before the new one: say the French for to regret, then the French for melancholy.
\end{quote}

\subsection*{You'll want to know: émouvoir}
\textbf{émouvoir} — \textquotedblleft{}to move, to touch (emotionally)\textquotedblright{}.

\begin{quote}
Sa lettre m'a beaucoup ému.
\end{quote}

\begin{quote}
\emph{Her letter moved me a lot.}
\end{quote}

Past participle \emph{ému}: \emph{j'étais ému} — I was moved.

\begin{cousinweb}
\emph{émouvoir} comes from Latin \emph{ēmovēre} \textquotedblleft{}to move out, to stir\textquotedblright{}, built on \emph{movēre} \textquotedblleft{}to move\textquotedblright{} (Lewis \& Short).
\end{cousinweb}

\begin{grammarlens}[title={What you've built}]
One more word for telling, explaining and describing at B1.
\end{grammarlens}

\subsection*{Your turn}
\begin{itemize}
\raggedright
  \item \emph{Say it:} \emph{émouvoir}
  \item \emph{Say it:} \emph{émouvoir}, once more
  \item \emph{Say it:} a sentence of your own with \emph{émouvoir}
  \item \emph{Recall:} say the French for to regret, then the French for melancholy, then say \emph{émouvoir} again
\end{itemize}

\begin{tcolorbox}[breakable,colback=teal!4,colframe=teal!35!black,title={Before you move on}]
What does émouvoir mean? (\textquotedblleft{}To move, to touch (emotionally)\textquotedblright{}.) Say it once more.
\end{tcolorbox}
\section[bouleverser]{bouleverser — to deeply upset, to overwhelm; to turn upside down}

\label{lesson:FR-C472-bouleverser}



\begin{quote}
Before the new one: say the French for melancholy, then the French for to move.
\end{quote}

\subsection*{You'll want to know: bouleverser}
\textbf{bouleverser} — \textquotedblleft{}to deeply upset, to overwhelm; to turn upside down\textquotedblright{}.

\begin{quote}
Cette nouvelle a bouleversé toute la famille.
\end{quote}

\begin{quote}
\emph{This news deeply upset the whole family.}
\end{quote}

Also of plans or a room: \emph{il a bouleversé nos projets}.

\begin{cousinweb}
\emph{bouleverser} is built from \emph{boule} \textquotedblleft{}ball\textquotedblright{}, from Latin \emph{bulla} \textquotedblleft{}a bubble\textquotedblright{}, and \emph{verser} \textquotedblleft{}to tip over\textquotedblright{}, from Latin \emph{versāre} \textquotedblleft{}to turn\textquotedblright{}, built on \emph{vertere} (Lewis \& Short): to send rolling like a ball.
\end{cousinweb}

\begin{grammarlens}[title={What you've built}]
One more word for telling, explaining and describing at B1.
\end{grammarlens}

\subsection*{Your turn}
\begin{itemize}
\raggedright
  \item \emph{Say it:} \emph{bouleverser}
  \item \emph{Say it:} \emph{bouleverser}, once more
  \item \emph{Say it:} a sentence of your own with \emph{bouleverser}
  \item \emph{Recall:} say the French for melancholy, then the French for to move, then say \emph{bouleverser} again
\end{itemize}

\begin{tcolorbox}[breakable,colback=teal!4,colframe=teal!35!black,title={Before you move on}]
What does bouleverser mean? (\textquotedblleft{}To deeply upset, to overwhelm; to turn upside down\textquotedblright{}.) Say it once more.
\end{tcolorbox}
